\documentclass[1pt]{article}
\usepackage[a4paper,margin=0.2in]{geometry} % Small margins
\usepackage{multicol} % For multiple columns
\usepackage{amsmath, amssymb} % For math symbols
\usepackage{bm} % Bold math
\usepackage{graphicx} % To insert small images
\usepackage{array}
\usepackage{booktabs}
\usepackage{placeins}
\usepackage{enumitem} % To make lists compact
\setlength{\columnsep}{1pt} % Reduce space between columns
\raggedcolumns{}
\begin{document}
\begin{multicols}{3} % 3 columns
\raggedcolumns
\setlength{\columnseprule}{0.1pt}
\scriptsize% Use a smaller font
\noindent % Avoid extra space at beginning
\abovedisplayskip=2pt
\belowdisplayskip=2pt


\textbf{ Order of Operations}
\begin{enumerate}[noitemsep, topsep=0pt]
    \item \textbf{Counting Degrees of Freedom:} 
holonomic constraint:
\begin{equation*}
    F\left(\vec{r}_1, \vec{r}_2 \ldots, \vec{r}_N, t\right)=0
\end{equation*}
number of free DOF \( n=3N-k \) where \( k \) is the number of constraints.
    \item{Find \( n \) generalized coordinates.}  
    \item{Express in terms of generalized coordinates}  
    \item{Calculating the Kinetic and Potential Energy}  
    \item {Writing the Equations of Motion}
    \end{enumerate}
\textbf{Generalized Kinetic Energy}
\begin{equation*}
    T = \frac{1}{2} m (\dot{x}^2 + \dot{y}^2 + \dot{z}^2) \quad \text{(Cartesian)}
\end{equation*}
\begin{equation*}
    T = \frac{1}{2} m (\dot{r}^2 + r^2 \dot{\theta}^2 + r^2 \sin^2 \theta \dot{\phi}^2) \quad \text{(Spherical)}
\end{equation*}
\begin{equation*}
    T = \frac{1}{2} m (\dot{r}^2 + r^2 \dot{\theta}^2 + \dot{z}^2) \quad \text{(Cylindrical)}
\end{equation*}

\textbf{ (M and K Matrices)}
\begin{equation*}
       T = \frac{1}{2} M_{ij} \dot{q}_i \dot{q}_j=\frac{1}{2} \mathbf{\dot{q}}^T \mathbf{M} \mathbf{\dot{q}}
\end{equation*}
\begin{equation*}
    V = \frac{1}{2} K_{ij} q_i q_j=\frac{1}{2} \mathbf{q}^T \mathbf{K} \mathbf{q}
\end{equation*}
\begin{equation*}
M_{ij} = \frac{\partial^2 \mathcal{L}}{\partial \dot{q}_i \partial \dot{q}_j} \ \quad
K_{ij} = -\frac{\partial^2 \mathcal{L}}{\partial q_i \partial q_j} 
\end{equation*}

\[
M_{ij} =
\begin{bmatrix}
\frac{\partial^2 \mathcal{L}}{\partial \dot{q}_1 \partial \dot{q}_1} & \frac{\partial^2 \mathcal{L}}{\partial \dot{q}_1 \partial \dot{q}_2} & \frac{\partial^2 \mathcal{L}}{\partial \dot{q}_1 \partial \dot{q}_3} \\
\frac{\partial^2 \mathcal{L}}{\partial \dot{q}_2 \partial \dot{q}_1} & \frac{\partial^2 \mathcal{L}}{\partial \dot{q}_2 \partial \dot{q}_2} & \frac{\partial^2 \mathcal{L}}{\partial \dot{q}_2 \partial \dot{q}_3} \\
\frac{\partial^2 \mathcal{L}}{\partial \dot{q}_3 \partial \dot{q}_1} & \frac{\partial^2 \mathcal{L}}{\partial \dot{q}_3 \partial \dot{q}_2} & \frac{\partial^2 \mathcal{L}}{\partial \dot{q}_3 \partial \dot{q}_3}
\end{bmatrix}
\]

\[
K_{ij} =
-\begin{bmatrix}
\frac{\partial^2 \mathcal{L}}{\partial q_1 \partial q_1} & \frac{\partial^2 \mathcal{L}}{\partial q_1 \partial q_2} & \frac{\partial^2 \mathcal{L}}{\partial q_1 \partial q_3} \\
\frac{\partial^2 \mathcal{L}}{\partial q_2 \partial q_1} & \frac{\partial^2 \mathcal{L}}{\partial q_2 \partial q_2} & \frac{\partial^2 \mathcal{L}}{\partial q_2 \partial q_3} \\
\frac{\partial^2 \mathcal{L}}{\partial q_3 \partial q_1} & \frac{\partial^2 \mathcal{L}}{\partial q_3 \partial q_2} & \frac{\partial^2 \mathcal{L}}{\partial q_3 \partial q_3}
\end{bmatrix}
\]


\begin{align*}
    M_{ij} \ddot{q}_j + K_{ij} q_j = 0
\end{align*}
\textbf{generalized velocity and acceleration}
\begin{equation*}
    \dot{r}_a = \sum_i \frac{\partial r_a}{\partial q_i} \dot{q}_i
\end{equation*}
\begin{equation*}
    \ddot{r}_a = \sum_i \frac{\partial r_a}{\partial q_i} \ddot{q}_i + \sum_{i,j} \frac{\partial^2 r_a}{\partial q_i \partial q_j} \dot{q}_i \dot{q}_j
\end{equation*}
\textbf{Generalized Force}
\begin{equation*}
    Q_{j}=\sum_{i}\vec{F}_{i}\frac{\partial r_{i}}{\partial q_{j}}=-\sum_i \vec{\nabla}_{r_i} V \frac{\partial \overrightarrow{r_i}}{\partial q_j}=-\frac{\partial V}{\partial q_j}
\end{equation*}
\begin{equation*}Q_j^{(n c)}=\sum_i \vec{F}_i^{(n c)} \cdot \frac{\partial \vec{r}_i}{\partial q_j}=\sum_i \vec{F}_i^{(n c)} \cdot \frac{\partial \vec{v}_i}{\partial \dot{q}_j}\end{equation*}

\textbf{conservative force}
\begin{equation*}
     \mathbf{F} = -\nabla V
\end{equation*}
\textbf{central force}
\begin{equation*}
    \mathbf{F} =-\vec{\nabla} V(r)=-\frac{\partial V}{\partial r} \hat{r}
\end{equation*}
\begin{equation*}
    W = \int \mathbf{F} \cdot d\mathbf{r}
\end{equation*}
\[\vec{J}=\intop_{t_{1}}^{t_{2}}Fdt\]

\textbf{Angular Momentum and Torque}
\begin{equation*}
    \mathbf{L} = \mathbf{r} \times \mathbf{p}=mr^2=I\omega
\end{equation*}
\begin{equation*}
    \tau= \mathbf{r} \times \mathbf{F} = \frac{d\mathbf{L}}{dt}=mr^{2}\dot{\omega}=\mathbf{I}\dot{\omega}
\end{equation*}

\textbf{Non-Inertial Systems}
\begin{equation*}
    \mathbf{F}_{\text{eff}} = \mathbf{F} - m \mathbf{a}_{\text{frame}} - 2m (\mathbf{\Omega} \times \mathbf{v}) - m \mathbf{\Omega} \times (\mathbf{\Omega} \times \mathbf{r})
\end{equation*}

\begin{equation*}
    \mathbf{F}_{\text{centrifugal}} = -m \mathbf{\Omega} \times (\mathbf{\Omega} \times \mathbf{r})
\end{equation*}

\begin{equation*}
    \mathbf{F}_{\text{Coriolis}} = -2m (\mathbf{\Omega} \times \mathbf{v})
\end{equation*}

\textbf{Lagrangian Mechanics}
 \begin{equation*}
    \frac{d}{dt} \left( \frac{\partial \mathcal{L}}{\partial \dot{q}} \right) - \frac{\mathcal{L}}{\partial q} = 0
\end{equation*}
\begin{align*}
    {\frac{\partial \mathcal{\mathcal{L}}}{\partial q}}-{\frac{d}{d t}}\left({\frac{\partial \mathcal{L}}{\partial{\dot{q}}}}\right)+{\frac{d^{2}}{d t^{2}}}\left({\frac{\partial \mathcal{L}}{\partial{\ddot{q}}}}\right)=0
\end{align*}

\textbf{with force}
\begin{equation*}
    \frac{d}{dt}\left[\frac{\partial T}{\partial\dot{q}_{j}}\right]-\frac{\partial T}{\partial q_{j}}=-\frac{\partial V}{\partial q_j}=Q_{j}
\end{equation*}
 \begin{equation*}
     \frac{d}{d t}\left[\frac{\partial \mathcal{L}}{\partial \dot{q}_j}\right]-\frac{\partial \mathcal{L}}{\partial q_j}=Q_j^{(n c)}=\sum_{i}\vec{F}_{i}\frac{\partial r_{i}}{\partial q_{j}}
 \end{equation*}
      \textbf{Rayleigh  Function:}
    \begin{equation*}
        \mathbf{F = \frac{1}{2} \sum_{i,j} C_{ij} \dot{q}_i \dot{q}_j}
    \end{equation*}
    \begin{equation*}
    F = \frac{1}{2} \sum_{i} C_i \dot{q}_i^2
\end{equation*}
    
     \textbf{Dissipative Force:}
    \begin{equation*}
        \mathbf{Q_i^{\text{(diss)}} = - \frac{\partial F}{\partial \dot{q}_i}}=- c \dot{x}
    \end{equation*}
    
     \textbf{Euler-Lagrange Equation with Dissipation:}
    \begin{equation*}
        \mathbf{\frac{d}{dt} \left( \frac{\partial \mathcal{L}}{\partial \dot{q}_i} \right) - \frac{\partial \mathcal{L}}{\partial q_i} = -\frac{\partial F}{\partial \dot{q}_i}}
    \end{equation*}
    
     \textbf{Power Dissipated:}
    \begin{equation*}
        \mathbf{P_{\text{diss}} = \sum_{i} Q_i^{\text{(diss)}} \dot{q}_i = 2F}
    \end{equation*}
    

\textbf{Cyclic Coordinates Conservation}
\begin{equation*}
    \frac{\partial \mathcal{L}}{\partial q_i} = 0 \implies \frac{d}{dt}\frac{\partial \mathcal{L}}{\partial \dot{q}_i} =\frac{d}{dt}p_i = 0 \implies\\ p_i=\text{const.}
\end{equation*}
\textbf{Effective Lagrangian}
\begin{equation*}
        \mathcal{L}_{\text{effective}}=\mathcal{L}-p_i \dot{q}_i\quad \texttt{(for cyclic coordinates)}
\end{equation*}
    
\textbf{Canonical momentum}
\begin{equation*}
    p_i = \frac{\partial \mathcal{L}}{\partial \dot{q_i}}
\end{equation*}
\begin{equation*}
    S[q(t)] = \int_{t_1}^{t_2} \mathcal{L}(\dot{q},q,t) \, dt
\end{equation*}
\textbf{Noether's Theorem}
\begin{equation*}
    \frac{d}{dt}\left(\sum_i \frac{\partial \mathcal{L}}{\partial \dot{q}_i} \frac{\partial f}{\partial \alpha}\right) = 0
\end{equation*}
\begin{equation*}
    \frac{\partial\mathcal{L}}{\partial t}=0\Rightarrow E=\text{const.}
\end{equation*}
\textbf{Hamiltonian Mechanics}
\begin{equation*}
    \mathcal{H} = \sum_i \dot{q}_i p_i - \mathcal{L}
\end{equation*}
\begin{equation*}
    p_i = \frac{\partial \mathcal{L}}{\partial \dot{q_i}}
\end{equation*}
\begin{equation*}
    \dot{q_i} = \frac{\partial \mathcal{H}}{\partial p_i}=\{\mathcal{H},q_i\}, \quad \dot{p_i} = -\frac{\partial \mathcal{H}}{\partial q_i}=-\{\mathcal{H},p_i\}
\end{equation*}
\[\frac{d}{dt} f \left( p_i, q_i, t \right) = \{ \mathcal{H}, f \} + \frac{\partial f}{\partial t}
\]
\[\frac{d}{dt} \{ f, g \} = \frac{\partial \{ f, g \}}{\partial t} + \{ \{ f, g \}, \mathcal{H} \}
\]

\textbf{Hamiltonian for Common Systems}
\begin{equation*}
   \mathcal{H} = \frac{p^2}{2m} + \frac{1}{2} k x^2 \quad \text{(Harmonic Oscillator)}
\end{equation*}
\begin{equation*}
    \mathcal{H} = \frac{p^2}{2m} - \frac{GMm}{r} \quad \text{(Kepler Problem)}
\end{equation*}

\textbf{Poisson Brackets}
\begin{equation*}
    \{ f, g \} = \sum_i \left( \frac{\partial f}{\partial q_i} \frac{\partial g}{\partial p_i} - \frac{\partial f}{\partial p_i} \frac{\partial g}{\partial q_i} \right)
\end{equation*}
\begin{align*}
    q_i, q_j = \{p_i, p_j\} = 0, \quad \{p_i, q_j\} = \delta_{ij}
\end{align*}
\begin{equation*}
    \{A, B\} = -\{B, A\}
\end{equation*}
\begin{equation*}
    \{A, BC\} = B\{A, C\} + \{A, B\}C
\end{equation*}
\begin{align*}
    \{{f, \{g,h\}\} + \{g, \{h,f\}\} + \{h, \{f,g\}\} = 0}
\end{align*}
\begin{align*}
   \{f(u), g(v)\} = \{u,v\} \frac{\partial f}{\partial u} \frac{\partial g}{\partial v}
\end{align*}
\begin{align*}
   \{ f(\vec{p}), g(\vec{q})\} = \sum \frac{\partial f}{\partial p_i} \frac{\partial g}{\partial q_i}
\end{align*}
\begin{align*}
    f(\vec{q}, \vec{p}), q_i\} = \frac{\partial f}{\partial p_i}, \quad \{ f(\vec{q}, \vec{p}), p_i\} = -\frac{\partial f}{\partial q_i}
\end{align*}
\begin{align*}
    \{L_i, L_j\} = -\sum_k \epsilon_{ijk} L_k
\end{align*}
\textbf{Constant of Motion}
\begin{equation*}
    \frac{dA}{dt} = \{ A, H \} + \frac{\partial A}{\partial t}
\end{equation*}


\textbf{Canonical Transformations Condition}
\begin{equation*}
    \sum_i (p_i dq_i - P_i dQ_i) = dF
\end{equation*}

\textbf{Generating Functions}
\begin{center}
\begin{tabular}{|c|l|}
\hline
\(F_1(q,Q)\) & \( p=\dfrac{\partial F_1}{\partial q},\quad P=-\dfrac{\partial F_1}{\partial Q}\) \\[4pt] \hline
\(F_2(q,P)\) & \( p=\dfrac{\partial F_2}{\partial q},\quad Q=\dfrac{\partial F_2}{\partial P}\) \\[4pt] \hline
\(F_3(p,Q)\) & \( q=-\dfrac{\partial F_3}{\partial p},\quad P=-\dfrac{\partial F_3}{\partial Q}\) \\[4pt] \hline
\(F_4(p,P)\) & \( q=-\dfrac{\partial F_4}{\partial p},\quad Q=\dfrac{\partial F_4}{\partial P}\) \\[4pt] \hline
\end{tabular}
\end{center}

\textbf{Damped Harmonic Oscillator}
\begin{equation*}
    \ddot{x} + 2\gamma\dot{x} + \omega_0^2 x = f(x)
\end{equation*}

\textbf{Fourier Transform Solution}
\begin{equation*}
    x_p = \frac{1}{\sqrt{2\pi}} \int_{-\infty}^{\infty} \frac{\hat{f}(\omega) e^{i\omega t} d\omega}{\omega_0^2 - \omega^2 + 2i\gamma\omega}
\end{equation*}

\textbf{Forced Oscillation}
\begin{align}
    x_p &= \frac{f_0 \cos(\omega_d t + \varphi)}{\sqrt{(\omega_0^2 - \omega_d^2)^2 + (2\gamma\omega_d)^2}} \\
    \varphi &= \tan^{-1} \left( \frac{2\gamma\omega_d}{\omega_0^2 - \omega_d^2} \right)
\end{align}

\textbf{Matrix Inverse}
\begin{equation*}
    A = \begin{pmatrix} a & b \\ c & d \end{pmatrix} \Rightarrow A^{-1} = \frac{1}{\det A} \begin{pmatrix} d & -b \\ -c & a \end{pmatrix}
\end{equation*}
\textbf{Small Oscillations} Solution Scheme

We assume that the system has an equilibrium at $\dot{\mathbf{q}}_0$. We define a new coordinate system $q_i$ based on a displacement from equilibrium. Expanding the Lagrangian to second order, we identify:
First-order terms in time vanish. Constants remain.

Expressed in the form:
\[
\mathcal{L}_{\text{eff}} \approx \frac{1}{2} \dot{\mathbf{q}}^T M \dot{\mathbf{q}} - \frac{1}{2} \mathbf{q}^T K \mathbf{q}
\]

Now, we solve the eigenvalue problem:
\[
\det \left( K - \omega^2 M \right) = 0
\]
The final solution $\vec{\omega}$ the eigenvalues vector with the $\vec{u_i}$ the eigenvector
\[\vec{q}(t)\approx\sum_i \vec{u_i}\cos(\omega_it)\]
\textbf{Oscillatory Motion Period}
\begin{align*}
        t = \int_{x_0}^{x} \frac{dx}{\sqrt{\frac{2}{m} (E - V(x))}} \\
        \tau = 2 \int_{x_1(E)}^{x_2(E)} \frac{dx}{\sqrt{\frac{2}{m} (E - V(x))}}
\end{align*}



 $x_1(E)$ and $x_2(E)$ are the roots of the equation $E = V(x)$, representing the turning points.

\textbf{Orbital Motion}
\begin{equation*}
    \frac{d^2u}{d\theta^2} + u = -\frac{m}{L^2} \frac{dV}{du}
\end{equation*}
\begin{equation*}
    r = \frac{a(1 - e^2)}{1 + e\cos\theta}
\end{equation*}
\begin{equation*}
    r(\theta) = \frac{r_0}{1 + \epsilon\cos(\theta-\theta_0)}=\frac{a(1-\epsilon^2)}{1 + \epsilon\cos(\theta-\theta_0)}
\end{equation*}
\begin{align*}
    r_0 = \frac{p_\theta^2}{\mu G m_1 m_2}=a(1-\epsilon ^2)
\end{align*}
\begin{align*}
    r_{\min} = \frac{r_0}{1 + \epsilon}=a(1-\varepsilon), \quad r_{\max} = \frac{r_0}{1 - \epsilon}=a(1+\varepsilon)
\end{align*}
\begin{align*}
    a=\frac{r_{max}+r_{min}}{2},\ b=a\sqrt{1-\varepsilon^2}
\end{align*}
\begin{align*}
    \varepsilon=\sqrt{1+\frac{2E p_\theta^{2}}{G^{2}\mu M_{1}^{2}M_{2}^{2}}}=\sqrt{1-\left(b/a\right)^{2}}
\end{align*}

\begin{equation*}
    E = \frac{1}{2} m v^2 - \frac{GMm}{r}
\end{equation*}
\begin{align*}
    \dot{r} = \sqrt{\frac{2}{\mu} \left( E - V(r) - \frac{\ell^2}{2\mu r^2} \right)}
\end{align*}
\begin{equation*}
    L = m r^2 \dot{\theta} = \text{constant}
\end{equation*}
\begin{align*}
    \varepsilon<1, \quad r(\theta) = \frac{r_0}{1 + e \cos(\theta - \theta_0)}
\end{align*}
\begin{align*}
    \varepsilon = 1, \quad r(\theta) = \frac{p_\theta^2}{2\mu G m_1 m_2}
\end{align*}
\begin{align*}
    \varepsilon>1, \quad r(\theta) = \frac{r_0}{1 + \varepsilon \cos(\theta - \theta_0)}
\end{align*}
\begin{center}
\begin{tabular}{|>{\raggedright\arraybackslash}p{0.05\linewidth}|>{\raggedright\arraybackslash}p{0.35\linewidth}|l|} \hline 

Orbit & Condition & Key Eq.\\[2pt] \hline 

Circ. & \(E=E_{\min},\,\epsilon=0\) & \(r_0=\dfrac{L^2}{\mu G m_1m_2}\)\\[4pt] \hline  
Ellip. & \(E_{\min}<E<0,\,0<\epsilon<1\) & \(r=\dfrac{Gm_1m_2}{2E}\Bigl(1\pm\sqrt{1+\frac{2EL^2}{\mu (Gm_1m_2)^2}}\Bigr)\)\\[4pt] \hline  
Parab. & \(E=0,\,\epsilon=1\) & \(r_{\text{par}}=\dfrac{L^2}{2\mu Gm_1m_2}\)\\[4pt] \hline 
Hyp. & \(E>0,\,\epsilon>1\) & \(r_{\text{hyp}}=\dfrac{Gm_1m_2}{2E}(1+\epsilon)\)\\ \hline

\end{tabular}
\end{center}
\textbf{Kepler Two-Body Problem}
\begin{align*}
    L = \frac{\mu}{2} \left( \dot{r}^2 + r^2 \dot{\phi}^2 \right) - V(r).
\end{align*}
\begin{align*}
    \mu \ddot{r} = -\frac{dV}{dr} + \frac{p_\theta^2}{\mu r^3}
\end{align*}
\begin{align*}
    \mu=\frac{m_1m_2}{m_1+m_2}
\end{align*}

\text{in the COM frame}
\begin{align*}
    L = \frac{\mu \dot{r}^2}{2} - V(r)
\end{align*}
\begin{align*}
    L = \mathbf{r} \times \mu \dot{\mathbf{r}} = \text{const}
\end{align*}
\begin{align*}
    P_{\phi} = \frac{\partial L}{\partial \dot{\phi}} = \mu r^2 \dot{\phi} = \text{const}
\end{align*}
\begin{align*}
      \mu \ddot{\mathbf{r}} = -\frac{GM\mu}{r^2} \hat{\mathbf{r}} = -\frac{\partial V_{\text{eff}}}{\partial r}
\end{align*}
    \begin{align*}
        V_{\text{eff}}(r) = V(r) + \frac{p_\theta^2}{2\mu r^2}=- \frac{G M_1 M_2}{r} + \frac{l^2}{2\mu r^2},
    \end{align*}
   
\begin{align*}
    E = \frac{1}{2} \mu \dot{r}^2 + V_{\text{eff}}(r)
\end{align*}
\begin{align*}
    \Delta\phi\equiv\mathrm{~{2}}\int_{r_{1}}^{r_{2}}\frac{L d r}{r^{2}\sqrt{2\mu(E-V_{e f f}(r))}}
\end{align*}


\textbf{Laplace-Runge-Lenz Vector}
\begin{equation*}
    \mathbf{A} = \mathbf{p} \times \mathbf{L} - m k \hat{\mathbf{r}}
\end{equation*}

\textbf{Rigid Body Dynamics}
\begin{equation*}
    \mathbf{R}_{CM} = \frac{1}{M} \sum_i m_i \mathbf{r}_i
\end{equation*}
\begin{equation*}
    I_{ij} = \sum_k m_k \left( r_k^2 \delta_{ij} - x_{k_i} x_{k_j} \right)
\end{equation*}
\begin{equation*}
    I = I_C + Md^2
\end{equation*}
\begin{equation*}
\begin{cases}
    I_1 \dot{\omega}_1 + (I_3 - I_2) \omega_2 \omega_3 = \tau_1 \\
    I_2 \dot{\omega}_2 + (I_1 - I_3) \omega_3 \omega_1 = \tau_2 \\
    I_3 \dot{\omega}_3 + (I_2 - I_1) \omega_1 \omega_2 = \tau_3
\end{cases}
\end{equation*}


\textbf{Moment of Inertia }
\begin{equation}
    I = \sum m_{\alpha} r_{\alpha}^2
\end{equation}
\[I = \int_V \rho(\mathbf{r})\, r^2\, dV\]
\begin{equation*}
    I_{xx} + I_{yy} = I_{zz}
\end{equation*}
\centering
\renewcommand{\arraystretch}{1.2}
\setlength{\tabcolsep}{3pt}
\renewcommand{\arraystretch}{1.1}
\begin{tabular}{|l |c|}
\toprule
\textbf{Object} & \textbf{\(I\)} \\
\midrule
Solid Sphere & \(\frac{2}{5}MR^2\)\\[2mm] \hline 
Hollow Sphere & \(\frac{2}{3}MR^2\)\\[2mm] \hline 
Solid Cylinder/Thin Disk & \(\frac{1}{2}MR^2\)\\[2mm] \hline 
Thin Rod (center) & \(\frac{1}{12}ML^2\)\\[2mm] \hline 
Rod (end) & \(\frac{1}{3}ML^2\)\\[2mm] \hline 
Rectangle (center) & \(\frac{1}{12}M(a^2+b^2)\)\\[2mm] \hline 
Ring (center) & \(MR^2\)\\[2mm] \hline 
Solid Cone (sym. axis) & \(\frac{3}{10}MR^2\)\\[2mm]
Cube (center) & \(\frac{1}{6}Ma^2\)\\
\bottomrule \hline
\end{tabular}

\[T=T_{C.M}+T_{Rot}=\frac{M}{2}V_{C.M}^2+\frac{I_{C.M}\omega^2}{2}=\]
\[=\frac{M}{2}V_{cm}^{2}+\frac{1}{2}\vec{\omega}^{T}I\vec{\omega}\]
\[T_{rot}=\frac{1}{2}\sum_{i,j}\omega_{i}I_{i,j}\omega_{j}\]

\textbf{Scattering }
For scattering with impact parameter $\rho$, scattering angle $\theta$, asymptotic velocity $v_\infty$, and reduced mass $\mu$, the differential cross section is:
\begin{equation*}
    \frac{d\sigma}{d\Omega}
    \;=\;
    \frac{\rho(\theta)}{\sin \theta}
    \left|\frac{d\rho}{d\theta}\right|.
\end{equation*}
The total cross section is:
\begin{equation*}
    \sigma_{T}
    \;=\;
    2\pi \int_{0}^{\pi}
    \rho(\theta)
    \left|\frac{d\rho}{d\theta}\right|
    \,d\theta=2\pi\intop_{0}^{\pi}\frac{d\sigma}{d\Omega}\sin\theta d\theta
\end{equation*}
\[\Delta \theta = 2 \int_{r_{\min}}^{\infty} \frac{l / r^2}{\sqrt{2\mu \left( E - V(r) - \frac{l^2}{2\mu r^2} \right)}} \, dr
\]
We identify
\[
E \;=\;\frac12\,\mu\,v_{\infty}^{2},
\quad
\ell \;=\;\mu\,\rho\,v_{\infty},
\]
and if $r_{\min}$ is the distance of closest approach in the scattering, then
\[\ \rho^2=\frac{l^2}{\mu^2v_{\infty}^{2}}=\frac{l^2}{2\mu E}\]
\[E=V\left(r_{min}\right)+\frac{l^{2}}{2\mu r_{min}^{2}}\]
\[
\rho^{2}
\;=\;
r_{\min}^{2}\Bigl(1 \;-\;\frac{V(r_{\min})}{E}\Bigr).
\]

\begin{equation*}
    \theta \approx - \frac{2b}{\mu v_\infty^2} \int_b^\infty \frac{V'(r) \, dr}{\sqrt{r^2 - b^2}}
\end{equation*}
\begin{equation*}
    b = \frac{L}{\mu v_\infty}
\end{equation*}
\textbf{Legendre Transform}
Given a function \( f(x) \), its Legendre transform is defined by:
\begin{equation*}
    \frac{df}{dx} = f'(x) = s(x)
\end{equation*}
Thus, we obtain a new function:
\begin{equation*}
    g(s) = -s\cdot x(s) + f(x(s))
\end{equation*}
\textbf{Inverse Transform}
We assume that the transform of \( g(s) \) satisfies:
\begin{equation*}
    V(x) = g(s(x)) + f(x)
\end{equation*}
Thus, we get:
\begin{equation*}
    -x(s) \cdot xu(x)
\end{equation*}
\textbf{Homogeneous Function of Degree \( k \)}
\begin{equation*}
    f(a x) = a^k f(x), \quad \forall a
\end{equation*}
\begin{equation*}
    E_\text{Kin}(a\vec{v})=\frac{1}{ 2}ma^2v^2=a^2E_\text{Kin}(\vec{v})
\end{equation*}
\begin{equation*}
    V(r)=g r^k 
\end{equation*}

\textbf{Euler’s Theorem for Homogeneous Functions}
\begin{equation*}
    \nabla f(x) \cdot x = k f(x)
\end{equation*}
Alternatively, in differential form:
\begin{equation*}
    x \cdot \frac{\partial f(x)}{\partial x} = k f(x)
\end{equation*}
\textbf{Virial Theorem}
\begin{equation*}
    2 \langle T \rangle + \langle V \rangle = 0
\end{equation*}
\begin{equation*}
    2<K>=-<\sum_j \frac{\partial \mathcal{L}}{\partial q_j} q_j>=<\sum_i \vec{r}_i \frac{\partial V}{\partial \vec{r}_i}>
\end{equation*}
\[2<T>=k_{deg }<V>\]
\textbf{Scaling}
\[x\rightarrow\alpha x,\ t\rightarrow\beta t\Rightarrow v\rightarrow\frac{\alpha}{\beta}v\]
\[\tilde{\mathcal{L}}=\left(\frac{\alpha}{\beta}\right)^{2}T-\alpha^{k}V\]
\[\left(\frac{\alpha}{\beta}\right)^{2}=\alpha^{k}\Rightarrow\beta=\alpha^{1-\frac{k}{2}}\]
\begin{equation*}
    \frac{v^{\prime}}{v}=\left(\frac{l^{\prime}}{l}\right)^{\frac{1}{2} k}, \frac{E^{\prime}}{E}=\left(\frac{l^{\prime}}{l}\right)^k,\left(\frac{l^{\prime}}{l}\right)^{1-\frac{1}{2}k}=\frac{t'}{t}
\end{equation*}
\textbf{Buckingham $\Pi$ Theorem (d-distance)}
\begin{align*}
        f(x_1, ..., x_n) = 0 \Rightarrow f(\Pi_1, ..., \Pi_{n-r}) = 0
\end{align*}
\begin{equation*}
    \tau \sim d^{1 - \frac{k}{2}}
\end{equation*}

\begin{equation*}
    v \sim \frac{d}{\tau} \sim \frac{d}{d^{1 - \frac{k}{2}}} \sim d^{k/2}
\end{equation*}

\begin{equation*}
    E \sim m v^2 \sim v^2 \sim d^k
\end{equation*}

\begin{equation*}
    L \sim \frac{d^2}{\tau} \sim d^{1 + \frac{k}{2}}
\end{equation*}
\textbf{Elliptic Coordinates}
\begin{equation*}
    x = a r \cos\theta, \quad y = b r \sin\theta
\end{equation*}
The area element in these coordinates is:
\begin{equation*}
    dA = a b r \, dr \, d\theta
\end{equation*}
where the integral over the area is:
\begin{equation*}
    \int_0^{2\pi} \int_0^1 dA
\end{equation*}
\textbf{Binet’s Equation}
Binet's equation describes the motion of a particle under a central force using the inverse radial distance \( u = \frac{1}{r} \):
\begin{equation*}
u''+u=-\frac{\mu }{p_\theta^{2}}\frac{1}{u^{2}}F(\frac{1}{u}) \end{equation*}
where:
- \( u = \frac{1}{r} \) is the reciprocal of the radial coordinate,
- \( F(1/u) \) is the central force as a function of \( u \),
\textbf{Solution for an Inverse-Square Law Force}
For a force of the form:
\begin{equation*}
    F(r) = -\frac{k}{r^2}
\end{equation*}
Substituting into Binet’s equation:
\begin{equation*}
    \frac{d^2 u}{d\theta^2} + u = \frac{k}{m h^2}
\end{equation*}
The general solution is:
\begin{equation*}
    u(\theta) = \frac{1}{r} = \frac{k}{m h^2} + C \cos(\theta - \theta_0)
\end{equation*}
This solution describes conic sections (ellipses, parabolas, or hyperbolas)
For gravitation 
\begin{align*}
    u'' + u = \frac{\mu G m_1 m_2}{p_\theta^2}
\end{align*}
The solution is 
\begin{align*}
    r(\theta) = \frac{r_0}{1 + e \cos(\theta - \theta_0)}
\end{align*}
where
\begin{align*}
    r_0 = \frac{p_\theta^2}{\mu G m_1 m_2}
\end{align*}
\textbf{Vector and Tensor Calculus}

\begin{equation*}
    \mathbf{a} \cdot \mathbf{b} = a_i b_i
\end{equation*}
\begin{equation*}
    ( \mathbf{a} \times \mathbf{b} )_i = \epsilon_{ijk} a_j b_k
\end{equation*}
\begin{equation*}
    (\nabla \times \mathbf{b})_i = \epsilon_{ijk} \partial_j b_k
\end{equation*}
\begin{equation*}
    (AB)_{ij} = A_{ik} B_{kj}
\end{equation*}
\begin{equation*}
    v_i \frac{1}{2} m_{jk} v_j v_k = m_{ik} v_k
\end{equation*}

\textbf{Index Notation Identity}
\begin{equation*}
    \epsilon_{ijk} \epsilon_{ilm} = \delta_{jl} \delta_{km} - \delta_{jm} \delta_{kl}
\end{equation*}

\textbf{Levi-Civita}
\begin{align*}
    \epsilon_{123}=1,\ \epsilon_{213}=-1
    ,\ \epsilon_{221}=0
\end{align*}
\begin{equation*}
    \epsilon_{ijk}= \epsilon_{jki} = -\epsilon_{jik}
\end{equation*}
\begin{equation*}
    \hat{e}_i \times \hat{e}_j=\epsilon_{i j k} \hat{e}_k
\end{equation*}
\textbf{complete derivative}
\begin{equation*}
    \frac{d X}{d t}=\frac{\partial X}{\partial t}+\frac{\partial X}{\partial q} \frac{d q}{d t} 
\end{equation*}
\begin{align*}
        \mathbf{r}  & = r(\cos\varphi,\ \sin\varphi) = r \hat{\mathbf{r}}\ , \\[1ex]
   \dot{\mathbf{r}} &= \dot{r}\hat{r}+r\dot{\theta}\hat{\theta} ,\\[1ex]
  \ddot{\mathbf{r}} &=(\ddot{r}-r\dot{\theta}^2)\hat{r}+(2\dot{r}\dot{\theta}+r\ddot{\theta})\hat{\theta}
\end{align*}

\textbf{Rotation About the $x$-Axis (Roll)}
\begin{equation*}
R_x(\gamma) =
\begin{bmatrix}
1 & 0 & 0 \\
0 & \cos\gamma & -\sin\gamma \\
0 & \sin\gamma & \cos\gamma
\end{bmatrix}
\end{equation*}

\textbf{Rotation About the $y$-Axis (Pitch)}
\begin{equation*}
R_y(\beta) =
\begin{bmatrix}
\cos\beta & 0 & \sin\beta \\
0 & 1 & 0 \\
-\sin\beta & 0 & \cos\beta
\end{bmatrix}
\end{equation*}

\textbf{Rotation About the $z$-Axis (Yaw)}
\begin{equation*}
R_z(\alpha) =
\begin{bmatrix}
\cos\alpha & -\sin\alpha & 0 \\
\sin\alpha & \cos\alpha & 0 \\
0 & 0 & 1
\end{bmatrix}
\end{equation*}
\[R(\alpha, \beta, \gamma) = R_z(\alpha) R_y(\beta) R_x(\gamma)
\]
\textbf{angular velocity matrix}
\begin{align*}
\Omega=
    \begin{pmatrix} \omega_1 \\ \omega_2 \\ \omega_3 \end{pmatrix}
=
\begin{pmatrix} 
\dot{\phi} \sin\theta \sin\psi + \dot{\theta} \cos\psi \\ 
\dot{\phi} \sin\theta \cos\psi + \dot{\theta} \sin\psi \\ 
\dot{\psi} + \dot{\phi} \cos\theta 
\end{pmatrix}
\end{align*}
\textbf{Parametric Resonance}
\[
\ddot{x} + \omega (t) x = 0
\]
where $\omega (t)=\omega (t+T)\implies x^*(t)=x(t+T)$.

\textbf{Conditions for Parametric Resonance}
We substitute the solution space of the motion and obtain a set of solutions. Due to the periodicity of the equation, we obtain a solution in the form:
\[
x_i (t + nT) = \mu_i^n x_i=\sum_j A_{ij}\tilde{x}(t)
\]
\[x_i(t)=\mu_i^{t/T}\Pi_i(t)    \]
\[\left|\mu_{1}\right|\neq\left|\mu_{2}\right|\implies\textbf{Parametric Resonance}\]
where the solution can be expressed as a product of periodic functions. This leads to the condition:
\[
\mu_1 \mu_2 = 1 \Rightarrow x_2 x_1 - x_1 x_2 = \text{const}
\]
Thus, parametric resonance occurs if $\mu_1 \neq 1$.

\textbf{Taylor}
\[
(1 + x)^\alpha
 = 1
 + \alpha x
 + \frac{\alpha(\alpha-1)}{2!}\,x^2
 + \cdots
\]
\[
f(x) = f(a)
  + f'(a)\,(x - a)
  + \frac{f''(a)}{2!}\,(x - a)^2
   + \cdots
\]
\[
\sin x
 = x
 - \frac{x^3}{3!}
 + \frac{x^5}{5!}+\cdots
\]
\[
\cos x
 = 1
 - \frac{x^2}{2!}
 + \frac{x^4}{4!}
 + \cdots\]\
\[
\frac{1}{1 - x}
 = 1
 + x
 + x^2
 + x^3+ \cdots
\]


\end{multicols}
\end{document}